\documentclass[10pt,a4paper]{amsart}
\usepackage[margin=19mm]{geometry}
\usepackage{amsmath,amssymb,mathtools}
\usepackage[scr=rsfs,cal=euler]{mathalfa}
\usepackage{xfrac}
\usepackage[unicode,hidelinks,bookmarks=false]{hyperref}
\newcommand{\ie}{i.e.\ }
\newcommand{\cf}{cf.\ }
\newcommand{\resp}{respectively\ }
\newcommand{\Subsection}{\subsection}
\providecommand{\nobreakdash}{\nobreak\hbox{-}}
\setlength{\emergencystretch}{2em}
\multlinegap0pt
\usepackage{bm}
\usepackage{bbm}
\usepackage{dsfont}
\usepackage{xspace}
\usepackage{cancel}
\usepackage{mathtools}
\usepackage{amsthm}
\makeatletter
\@ifundefined{thm}{\newtheorem{thm}{Thm}}{}
\@ifundefined{prop}{\newtheorem{prop}[thm]{\bf Proposition}}{}
\makeatother
\def \L{\mathcal{L}}
\def \O{\mathcal{O}}
\def \Pic{\mathrm{Pic}}
\title[Hierarchical filtrations of line bundles and optimal algebra]{Hierarchical filtrations of line bundles and optimal algebraic geometry codes}
\author{Rahim Rahmati-asghar}
\date{Source: arXiv:2507.01859v6 (CC BY 4.0)}
\begin{document}
\maketitle

\noindent\emph{Source and licence.} Hierarchical filtrations of line bundles and optimal algebraic geometry codes. Version arXiv:2507.01859v6 (CC BY 4.0), distributed under the Creative Commons Attribution 4.0 International licence, as recorded in the arXiv licence metadata for this exact version and shown on the abstract page https://arxiv.org/abs/2507.01859v6. Full rights notice: licenses/arxiv-2507.01859-CC-BY-4.0.txt.\par\smallskip

\noindent\textit{Editorial scope.} Counted unit: the Prop whose source label is \texttt{prop:max-filtration} in the article (source0.tex lines 194--200, source label \texttt{prop:max-filtration}), delivered together with the article's own complete original proof of it (source0.tex lines 202--247, one \texttt{proof} environment of 46 lines). No mathematics is written, repaired or re-derived: statement and proof are the article's own text line for line. Required context retained uncounted: none. Every definition, display and label of those lines is retained unchanged. Omitted from this package and not redistributed: 1--193, 201--201, 248--1660. Nothing inside a retained excerpt is omitted or altered: every retained body is delivered exactly as the article's lines are written, so no omission receipt of any kind accompanies this package. Citations: the retained excerpts carry no citation command at all, so no bibliography entry of the article is transcribed and no reference is redistributed. Reference adaptation: none was needed and none was made; every reference inside the delivered unit resolves inside it, so no unresolved reference or citation is produced. Printed slips kept literal: no printed word, symbol, hypothesis or number of the counted statement or of its proof was altered. Layout and class: the article's own class is \texttt{amsart} with packages including amsmath, amssymb, amsfonts, enumerate, amsthm, amscd, babel, hyperref, paralist, verbatim, algorithm, algorithmicx, algpseudocode, color; the wrapper is the lane's pinned \texttt{amsart} editorial wrapper at 10pt on A4, which restates the theorem-like environments the retained text needs and adds the article's own macro definitions that the retained text needs (\texttt{\textbackslash L}, \texttt{\textbackslash O}, \texttt{\textbackslash Pic}), with extra wrapper packages (bm, bbm, dsfont, xspace, cancel, mathtools). The delivered numbering is the unit's own. Assets: no retained span contains a figure environment, a graphics-inclusion command, a PDF-page inclusion, a table or a TikZ picture; no image file is copied into this package, the package declares no asset dependency and no image file is redistributed with it. Licence: version arXiv:2507.01859v6 of this article is distributed under the Creative Commons Attribution 4.0 International licence, as recorded in the arXiv licence metadata for this exact version and shown on the abstract page https://arxiv.org/abs/2507.01859v6. A full rights notice is delivered at licenses/arxiv-2507.01859-CC-BY-4.0.txt. Characters: every retained excerpt is pure ASCII, so the delivered engine, XeLaTeX with its default Latin Modern text font, reports no missing character.
\begin{prop}\label{prop:max-filtration}
Let $X$ be a smooth projective variety over a field, and let $L$ be a line bundle on $X$.
If there exists at least one hierarchical filtration of $L$, then the set of possible
lengths of such filtrations is bounded above and hence admits a maximum. 
In particular, the hierarchical depth $h(L)$ is finite, and there exists a hierarchical
filtration of $L$ of maximal length.
\end{prop}

\begin{proof}
Fix an ample line bundle $\O_X(1)$ on $X$, and let $H=c_1(\O_X(1))\in \Pic(X)$. 
Let 
$$
\O_X=\L_0 \subset \L_1 \subset \cdots \subset \L_h=L
$$
be a hierarchical filtration of $L$. 
By definition, for each $i$ there exists a nonzero effective Cartier divisor $E_i$ on $X$ 
such that 
$$
\L_i \cong \L_{i-1}\otimes \O_X(E_i).
$$
Equivalently, there is a nonzero section 
$s_i \in H^0\!\big(X,\L_i\otimes\L_{i-1}^{-1}\big)$ 
whose divisor of zeros is precisely $E_i$. 

Multiplying the sections $s_1,\dots,s_h$ yields a nonzero section 
$$
s := s_1 \cdots s_h \in H^0(X,L),
$$
and its zero divisor satisfies
$$
\operatorname{div}(s) = E_1 + \cdots + E_h.
$$
Now intersecting with $H^{\dim X - 1}$ gives
$$
\sum_{i=1}^h (E_i \cdot H^{\dim X - 1})
= \operatorname{div}(s)\cdot H^{\dim X - 1}.
$$
Since $\operatorname{div}(s)$ is linearly equivalent to $c_1(L)$, the right-hand side depends
only on $L$, not on the chosen filtration. Set
$$
N := c_1(L)\cdot H^{\dim X - 1} \in \mathbb{Z}_{\ge 0}.
$$
Because $X$ is projective and $H$ is ample, one has 
$E\cdot H^{\dim X-1} > 0$ for every nonzero effective Cartier divisor $E$ on $X$. 
Hence each summand $E_i\cdot H^{\dim X-1}$ is at least $1$, so
$$
h \le \sum_{i=1}^h (E_i \cdot H^{\dim X - 1}) = N.
$$
Thus the length $h$ of any hierarchical filtration of $L$ is bounded above by $N$.
Consequently the set of all possible lengths is a nonempty finite subset of 
$\{0,1,2,\dots,N\}$, and therefore admits a maximum. 
By definition this maximum is $h(L)$, and any filtration realizing it is a 
filtration of maximal length.
\end{proof}

\end{document}
